\section{Origine del codice}
\label{sec:code-origin}

Per trasparenza si dichiara la provenienza di ogni componente. Il dettaglio
completo \`e in \code{docs/origine\_codice.md}; qui si riporta la sintesi.

\begin{table}[H]
\centering
\small
\caption{Origine dei moduli. \emph{Adattato} = derivato dal materiale del
corso; \emph{Nuovo} = scritto da zero per questo progetto;
\emph{Standard} = replica di architetture note (citate).}
\label{tab:origin}
\begin{tabular}{p{5.2cm}p{2.6cm}p{5.2cm}}
\toprule
Componente & Origine & Fonte / note \\
\midrule
Q\hbox{-}Learning, SARSA & Adattato &
  esercitazioni del corso (\code{control.py}); adattato all'API Gymnasium e
  agli stati discretizzati \\
MC prediction, TD(0) & Adattato &
  \code{prediction.py}, esercizio ``random walk'' \\
Policy $\varepsilon$-greedy & Adattato &
  \code{Esercitazione\_RL} \\
Discretizzatore & Nuovo &
  gli ambienti del corso erano gi\`a discreti \\
Preprocessing Atari & Adattato &
  \code{Tutorial\_Gymnasium}; wrapper Gymnasium (non SB3) \\
Replay buffer & Nuovo &
  concetto da Mnih~et~al.~\cite{mnih2015} \\
Rete Nature-CNN / MLP & Standard &
  architettura da~\cite{mnih2015}, implementazione PyTorch originale \\
Loop di training DQN & Nuovo &
  parte pi\`u originale; il corso forniva solo DQN basato su SB3 (vietato) \\
Double DQN & Nuovo &
  da~\cite{vanhasselt2016}, $\sim$5 righe sul DQN base \\
Runner multi-seed & Ispirato &
  pattern \code{run\_bias\_experiment} del corso \\
Plotting (bande $\pm\sigma$) & Nuovo (parz.\ adattato) &
  \code{moving\_average} dal corso; curve multi-seed nuove \\
Script per ambiente, tabelle & Nuovo &
  \code{scripts/} \\
\bottomrule
\end{tabular}
\end{table}

In sintesi: i \emph{quattro algoritmi tabellari} sono adattati dalle
esercitazioni; il \emph{DQN}, il discretizzatore, l'infrastruttura
multi-seed e tutti gli script sono \emph{scritti da zero}. Nessun codice
Stable-Baselines3 \`e usato nel core, come richiesto dalla consegna.
